
en
vertu de 4.5, l'implication pour \(X\) en r\'esulte. Supposons donc dor\'enavant \(X\) projectif, muni d'un faisceau ample \(H\).

Soient \(X_i\) les composantes irr\'eductibles de \(X\), munies de la structure r\'eduite induite. Appliquant 4.5 et EGA III 2.6.2 au morphisme fini et surjectif \(\coprod X_i\to X\), on ram\`ene la d\'emonstration que (iii) implique (i) et (iv) au cas o\`u \(X\) est int\`egre.

Si \(X\) est irr\'eductible, il est \'evident que (ii) implique (v), que (iii) implique (vi), et compte tenu de 2.1, que (v) implique (vi), et que (vi) pour \(X\) est \'equivalent \`a (vi) pour \(X_{\red}\) (X 4.4.1).

Il reste donc \`a voir que (vi) implique (i) et (iv), et que (iv) implique (i), en supposant \(X\) int\`egre. Mais alors, (vi) (resp. (iv)) implique que la famille \(\mathcal L\) des classes des \(L^{\otimes m}\) (\(m\in\mathbb N\)) satisfait \`a 3.13 (iii) (resp. (iv)), donc est limit\'ee. Il existe donc un faisceau inversible \(L\) sur un \(X_T\), avec \(T\) de type fini sur \(k\), qui limite \(\mathcal L\). Soient \(p,q\) deux entiers diff\'erents tels que \(L^{\otimes r}\) et \(L^{\otimes q}\) correspondent \`a la m\^eme composante connexe de \(T\), alors \(L^{\otimes(p-q)}\) est donc alg\'ebriquement \'equivalent \`a \(0\)\,; d'o\`u (vi) (resp. (iv)) implique (i).

Enfin, on a de m\^eme que (vi) implique qu'il existe un entier \(N\) tel que pour tout entier \(m\), \(L^{\otimes m}\otimes H^{\otimes N}\) soit ample. Alors pour tout \(m\), \((L^{\otimes m}\otimes H)^{\otimes N}\) est ample\,; d'o\`u que (vi) implique (iv), ce qui ach\`eve la d\'emonstration.

\textbf{Th\'eor\`eme 4.7.} \emph{Soient \(S\) un sch\'ema noeth\'erien, et \(X\) un \(S\)-sch\'ema propre. Alors}~:
\begin{enumerate}[label=(\roman*)]
\item \emph{Le morphisme d'inclusion \(\underline{\Pic}^\tau_{X/S}\to\underline{\Pic}_{X/S}\) est repr\'esentable par une immersion ouverte.}
\item
\emph{Si \(X\) est projectif et int\`egre sur \(S\), alors \(\underline{\Pic}^\tau_{X/S}\to\underline{\Pic}_{X/S}\) 
